% Chapter Template

\chapter{Conclusiones} % Main chapter title

Este capítulo presenta un resumen de los logros y las dificultades para llegar al estado actual de la solución y sugiere una serie de siguientes pasos para la evolución del sistema.

\label{Chapter5} % Change X to a consecutive number; for referencing this chapter elsewhere, use \ref{ChapterX}

\section{Conclusiones generales}

En este trabajo se propuso diseñar e implementar un sistema capaz de predecir características musicales a partir de audio de canciones. El objetivo principal fue automatizar
un proceso tradicionalmente manual y laborioso de clasificación musical y reducir la dependencia de servicios de tercecros para evitar costos recurrentes.
El trabajo está en producción, clasifica las canciones que ingresan al sistema, es funcional en el contexto de una empresa en crecimiento. 

En relación a los objetivos específicos, se puede concluir que:

%----------------------------------------------------------------------------------------

%----------------------------------------------------------------------------------------
%	SECTION 1
%----------------------------------------------------------------------------------------

\begin{itemize}
\item Se cumplieron los requerimientos técnicos al lograr la caracterización de las variables solicitadas a partir de los datos de Plusyc.
\item La extracción de variables acústicas DSP sirvió para las predicciones, y la inclusión de embeddings y probabilidades de PANNs fue lo que permitió resolver las variables semánticas complejas. Sin esto último, no se habría completado el trabajo.
\item Se aplicaron acciones para mitigar problemas en los datos. Fue necesario aplicar técnicas de preprocesamiento; se cortó el histórico en una fecha específica para aislar un cambio en la distribución de los datos.
\item Se entrenaron modelos de aprendizaje automático para predecir las características musicales y la evaluación de las métricas de rendimiento dan resultados productivos, acorde al análisis y diseño realizados.
\item La decisión de usar bibliotecas como Librosa y PANNs inference y modelos como redes MLP y ensambles de árboles, permitió no depender de licencias ni pagar costos por consultas a modelos de terceros, también facilitó el trabajo al no requerir de arquitecturas complejas.
\item En específico, la clasificación de estados de ánimo y géneros sugiere que la ambigüedad humana al etiquetar afecta el aprendizaje. Estos resultados indican que el equipo de curaduría necesita estandarizar los criterios para la ingesta manual de datos.
\item Para la detección de público en vivo y habla en canciones, se requirió de limpieza de ruido por falsas etiquetas apoyada en un modelo preentrenado con el conjunto de datos AudioSet, creado para la generación automática de subtítulos como [Aplausos] o [Rapping].
\item Completar el trabajo tomó más tiempo del planificado en el cronograma original. Hubo una subestimación del esfuerzo, dado que la cantidad de modelos y el proceso de aprendizaje requirieron más iteraciones de las previstas.

\end{itemize}

\section{Próximos pasos}

Los diseños actuales quedan disponibles y se pueden usar para iteraciones futuras. Algunos aspectos que se podrían mejorar o explorar son:

\begin{itemize}
\item El uso de las inferencias de PANNs y la exploración de otros modelos preentrenados, para mejorar la limpieza de algunos datos y corregir audios mal etiquetados.
Como ejemplos, el \textit{speechiness}, \textit{liveness} y la predicción de \textit{instrumentalness}: el modelo actual logra diferenciar una canción instrumental de una que no lo es, pero los valores en el medio no son claros. 
Se sugiere revisar puntualmente, ya que presenta ruido en los valores intermedios por falta de definiciones claras.
\item También se sugiere iterar sobre el \textit{tempo} para buscar una solución más robusta que el algoritmo de DSP.
\item Se propone empezar a usar el estado de ánimo secundario como filtro en Plusyc para mejorar la precisión de las recomendaciones musicales.
\item El sistema se puede ajustar para el despliegue bajo el modelo de software como servicio y así permitir ofrecer el servicio al público interesado.
\item El equipo de curadores sigue clasificando música. Se contempla automatizar un reentrenamiento periódico para que el sistema aprenda nuevos géneros cada vez que una clase alcance un mínimo de 100 muestras, con el uso de tecnologías como Airflow y MLflow para optimizar estos procesos, hoy manuales.
\item Se podría evaluar a futuro el uso de arquitecturas basadas en transformadores, si se cuenta con mayor capacidad de cómputo, para mejorar la clasificación de clases densas. 
También, complementar el sistema con modelos de aprendizaje profundo enfocados exclusivamente en música.
\end{itemize}
%----------------------------------------------------------------------------------------

%   SECTION 2

%----------------------------------------------------------------------------------------

% \section{Próximos pasos}



% Acá se indica cómo se podría continuar el trabajo más adelante.



% \section{Conclusiones generales}

% % Este capítulo resume los hallazgos del proyecto y establece la hoja de ruta para su evolución en Plusyc.

% \begin{itemize}
% \item Se cumplió la totalidad de los requerimientos técnicos al lograr la caracterización de las variables acústicas solicitadas.
% \item La ejecución tomó más tiempo del planificado en el cronograma original. Hubo una subestimación del esfuerzo, ya que la cantidad de modelos obligó a iterar más de lo previsto para encontrar resultados aceptables.
% \item Se aplicaron acciones para mitigar problemas en los datos. Se cortó el histórico en enero de 2026 para aislar un cambio en la distribución de los registros. Además, se usó submuestreo y limpieza con modelos preentrenados para compensar el desbalance en clases como liveness y speechiness.
% \item La extracción de variables acústicas tradicionales sirvió como base, pero la inclusión de embeddings y probabilidades de PANNs fue lo que permitió resolver las variables semánticas complejas. Sin esto último, no se habría completado el trabajo.
% \item La clasificación de estados de ánimo y géneros demostró que la ambigüedad humana al etiquetar afecta el aprendizaje. Estos resultados indican que el equipo de curaduría necesita estandarizar los criterios para la ingesta manual de datos.
% \item Se optó por usar PyTorch y PANNs para no depender de licencias ni pagar costos por consultas a modelos de terceros. El sistema se desplegó bajo el modelo de software como servicio para facilitar un procesamiento escalable y asíncrono.
% \end{itemize}

% \section{Próximos pasos}

% % Acá se indica cómo se podría continuar el trabajo más adelante.

% \begin{itemize}
% \item Los modelos actuales quedan como base para iteraciones futuras. Se debe revisar puntualmente instrumentalness, ya que presenta mucho ruido por falta de definiciones claras en los valores intermedios.
% \item Se deja como trabajo pendiente usar las inferencias de PANNs para automatizar la limpieza de datos en el histórico y corregir audios mal etiquetados.
% \item El equipo de curadores sigue clasificando música. Se contempla automatizar un reentrenamiento periódico para que el sistema aprenda nuevos géneros cada vez que una clase alcance un mínimo de 100 muestras.
% \item Se propone empezar a usar el estado de ánimo secundario como filtro en Plusyc para mejorar la precisión de las recomendaciones musicales.
% \item Se sugiere evaluar a futuro el uso de arquitecturas basadas en transformadores para audios, si se cuenta con mayor capacidad de cómputo, para buscar mejoras en la clasificación de clases densas como los géneros.
% \end{itemize}